\section{A Five-Axis Taxonomy}
\label{sec:taxonomy}

% 本节对应 slides/outline.md 的 P4 + P5。
% 图表: 五轴速览表 + TikZ 分类树 + 分布表。

\subsection{Design principles}

We require three properties of our axes:

\begin{enumerate}
    \item \textbf{Orthogonality.} Each paper is assignable on each axis
    independently, without consulting its other assignments. If assigning on
    one axis requires first deciding another, the axes are not independent.

    \item \textbf{Descriptive coverage.} An axis must accommodate at least
    80\% of the surveyed papers without a residual \code{other} category. An
    axis requiring an \code{other} bucket is measuring the wrong thing.

    \item \textbf{Diagnostic value.} The axis must expose structure --- trends,
    gaps, or concentrations of controversy. An axis that partitions the
    literature arbitrarily has no analytic value.
\end{enumerate}

\subsection{The axes}

\begin{table}[H]
\centering
\caption{The five axes. Axis E is the one that bears directly on the definition
of RSI.}
\label{tab:axes}
\begin{tabularx}{\textwidth}{@{}l p{0.34\textwidth} X@{}}
\toprule
\thead{Axis} & \thead{Question} & \thead{Values} \\
\midrule
A: Target
  & What component of the harness is modified?
  & \code{prompt}, \code{control-flow}, \code{code}, \code{tool},
    \code{memory}, \code{evaluator}, \code{context} \\

B: Signal
  & What licenses the system to consider a change an improvement?
  & \code{self-preference}, \code{external-verifier}, \code{human},
    \code{failure-trajectory} \\

C: Operator
  & How is the search space explored?
  & \code{refine}, \code{evolutionary}, \code{textual-gradient},
    \code{tree-search}, \code{meta-optimizer} \\

D: Guarantee
  & What backs the claim that a change is beneficial?
  & \code{none}, \code{empirical}, \code{formal},
    \code{formal-safety} \\

E: Depth
  & Does the improved system participate in the next improvement?
  & \code{one-shot}, \code{iterative}, \code{full-rsi},
    \code{co-evolution} \\
\bottomrule
\end{tabularx}
\end{table}

\subsection{Axis B and the problem of the verifier}

Axis B carries analytic weight that the others do not. When the system that
proposes a change is also the system that certifies it, the verification signal
is not independent of the thing being optimized. This is the structure of a
reward-hacking loop, and it is why \code{self-preference} carries a
distinctly different epistemic status from \code{external-verifier}.

We treat this as the organizing question of Section~\ref{sec:open}.

\subsection{Axis E and the boundary of recursion}

Axis E follows directly from the definition in
Section~\ref{sec:background}. The distinction between \code{iterative} and
\code{full-rsi} is whether the \emph{improver} is itself improved, not merely
its output.

This axis requires attribution discipline. \cite{zelikman2023stop}, whose title
invokes recursive self-improvement, states in its own text that because the
language model is not altered, the result is not full recursive self-improvement.
We record that qualification rather than overriding it, since the
disambiguation of this literature is part of what this review contributes.

By contrast, \cite{zhang2025dgm} explicitly claims to improve its own ability
to modify its own codebase, and is therefore assigned \code{full-rsi}.

\subsection{Resulting classification}

\begin{figure}[H]
\centering
% TODO: 用 TikZ 绘制分类树，根节点为 A 轴的七个取值
% preamble.tex 已加载 tikz + trees 宏包。
% 要求：一页能讲清。画不完说明分类还没想清楚。
\todo{TikZ classification tree, root = the deep-read corpus, children = A
axis values. Use the tikz trees library already loaded in preamble.tex.}
\end{figure}

\begin{table}[H]
\centering
\caption{Assignment of the deep-read corpus to the five axes. A dash denotes a
cell not yet determined --- see Section~\ref{sec:evidence} for read status.}
\label{tab:assignment}
\footnotesize
\begin{tabularx}{\textwidth}{@{}l l X l l@{}}
\toprule
\thead{ID} & \thead{Title} & \thead{A: Target} & \thead{C: Operator} & \thead{E: Depth} \\
\midrule
\paperID{arxiv-2310.02304} & STOP \cite{zelikman2023stop}
  & \code{code} & \code{refine}, \code{evolutionary} & \code{one-shot} \\
\paperID{arxiv-2408.08435} & ADAS \cite{hu2024adas}
  & \code{control-flow}, \code{code} & \code{tree-search} & \code{one-shot} \\
\paperID{arxiv-2504.15228} & SICA \cite{robeyns2025sica}
  & \code{code} & \code{evolutionary} & \code{iterative} \\
\paperID{arxiv-2505.22954} & DGM \cite{zhang2025dgm}
  & \code{code}, \code{evaluator} & \code{evolutionary} & \textbf{\code{full-rsi}} \\
\paperID{arxiv-2506.13131} & AlphaEvolve \cite{novikov2025alphaevolve}
  & \code{code} & \code{evolutionary} & \code{one-shot} \\
\paperID{arxiv-2608.07645} & MGM \cite{liu2026mgm}
  & \code{code} & \code{evolutionary} & \code{iterative} \\
\paperID{arxiv-2606.05922} & RHO \cite{pan2026rho}
  & \code{prompt}, \code{control-flow}, \code{tool}
  & \na & \code{iterative} \\
\paperID{arxiv-2603.25111} & SEVerA \cite{banerjee2026severa}
  & \code{code}, \code{control-flow} & \na & \na \\
\paperID{arxiv-2607.15524} & Recursive Harness SI \cite{lee2026recursiveharness}
  & \code{context} & \na & \code{co-evolution} \\
\bottomrule
\end{tabularx}
\end{table}

\subsection{What the distribution shows}

\todo{Answer in two or three paragraphs: (1) which cells are dense, i.e. what
the mainstream approach is; (2) which cells are empty, i.e. where the research
gaps are; (3) note that \code{full-rsi} on axis E has a single member, and
argue carefully whether this reflects genuine difficulty or merely low
interest. The sample is small --- do not over-read it.}

\subsection{Limitations of the taxonomy}

\todo{State explicitly where this scheme fails. At minimum acknowledge: the
sample is small (n = \todo{N}), so distributional claims must be cautious;
multi-valued assignment on axis A involves judgement, and the conflicts are
themselves diagnostic; axis C deliberately excludes the prompt-optimization
lineage, which is a trade-off rather than an oversight.}

% ----------------------------------------------------------------------
% 候选引用（核对后启用）:
%   \cite{ning2026codeharness}       Code as Agent Harness
%   \cite{zhang2026selfharness}      Self-Harness
%   \cite{lin2026notbenefit}         争议: 能改 != 有收益
% ----------------------------------------------------------------------
